\chapter{Conclusion and Outlook}\label{ch:conclusion}

\section{Answering the Research Questions}\label{sec:answers}

\paragraph{RQ0 (Main question).} By how much active, measurement-driven cache-aware decisions
outperform their passive, statically dimensioned counterparts can be quantified conclusively only
after the end-to-end runs per platform (Section~\ref{sec:outlook}). What is prepared is the
apparatus that isolates exactly this difference: the same axis-constant structure is permuted once
actively (measurement-driven layout, prefetch, and value-placement decision) and once passively
(statically dimensioned) and compared under an identical compiler and flag basis on hybrid CPUs and
Sapphire Rapids.

\paragraph{RQ1 (Composition).} The nineteen permutation axes (search-algorithm,
cache-traversal, mapping, path-compression, node-type, memory-layout, allocator, prefetch,
concurrency, serialization, telemetry, value-handle, ISA, index-organization, I/O-dispatch,
migration, filter, and the two queuing axes) span a configuration space of far more than $10^{11}$
permutations. State-of-the-art designs are reconstructed unambiguously through algorithm-profile
XMLs: 30 SOTA profiles (8 Rank-1 + 22 Rank-2/3) demonstrate that the axis decomposition carries
P01--P33. The variadic API specialization (1/2/$N$ parameters onto
\texttt{std::vector}/\texttt{std::map}/tuple) makes both algorithm- and container-oriented patterns
expressible on one engine. The non-fully-orthogonal axis couplings (node type $\times$ SIMD $\times$
path compression) are marked as such and explicitly restricted during permutation.

\paragraph{RQ2 (Measurement methodology).} Measurement is split into a production mode
(\texttt{COMDARE\_\allowbreak EXPERIMENT\_\allowbreak MODE=OFF}, no overhead) and an experiment mode
(the result aggregator is an integral execution-engine member). Profile-aware workload routing
selects the workload per permutation from the traversal tag, ensuring fair comparison between
range-query- and point-lookup-oriented designs. Each series is collected at three granularities ---
micro-, macro-, and whole-system benchmarking ---, making both individual design components and
whole algorithms comparable. Contradictory findings in the literature on the optimal node/cache-line
size (one versus up to sixteen cache lines) are thus not adopted from incompatible original setups
but re-measured axis-isolated, with all else held equal.

\paragraph{RQ3 (Probe comparison).} The concrete comparison of PRT-ART against the 30 SOTA
profiles is pending the first measurement runs (Section~\ref{sec:outlook}). The architecture is
fully prepared: the PRT-ART profile, the three mandatory series templates, the
\texttt{ExperimentDriver} with auto-pickup of the SOTA profiles, and the ABI-conformant adapter with
all \texttt{std::map} operations and notify hooks. The cheapest local page representation per branch
point (point versus compact page) and its platform-calibrated switching rules are likewise part of
the PRT-ART profile.

\paragraph{RQ4 (Compile time vs.\ run time).} The separation is anchored in the build system:
platform-specific axes (ISA, hardware constants) are fixed at compile time and produce
provenance-verifiable binaries per permutation (SHA256-validated original bodies,
\texttt{experiment\_compiler()} identifier), while layout and page selection remain selectable at
run time via a bit-encoded permutation identifier, without violating the common compiler and flag
basis.

\section{Limitations}\label{sec:limitations}

\begin{itemize}
  \item \textbf{Module bodies are currently stubs:} the code-generation pipeline produces
  ABI-conformant skeleton modules per profile; the actual per-permutation algorithm body remains for
  a follow-up phase.
  \item \textbf{Hardware validation pending:} configuration is green locally, but the full
  \texttt{ctest} and end-to-end \texttt{messung\_driver} runs across the BlockAO platforms
  (AVX2/AVX-512, ARM NEON/SVE) are outstanding.
  \item \textbf{Results chapter:} currently methodological; concrete measurements follow.
\end{itemize}

\section{Outlook}\label{sec:outlook}

In the short term: the end-to-end run of the three mandatory series on real hardware, the generation
of the measurement diagrams, and the completion of the per-permutation module bodies. In the medium
to long term: a GPU adapter (Grace Hopper cluster, ZIH), PIM-allocator integration (A16), and formal
verification of selected permutation classes (StarMalloc-style, A13). The three-repository
architecture lets future students add further probe algorithms analogously to PRT-ART without
modifying the cache-engine tool library --- a key design principle for the scalability of the
research infrastructure. Beyond the tree-like search structures treated here, the anatomy principle
extends through higher-level abstractions to the further structure classes of
Section~\ref{sec:search-classes} --- hash-based, spatial, as well as set- and sequence-oriented
structures ---, so that the same measurement and permutation apparatus can be applied to entirely
new classes without rebuilding it; herein lies a key lever for future experimental progress.
